%% kapa: exam=1 round=34 year=2023 period=2 form=A subject=appraisal_law subject_ko=감정평가관계법규 seq=28 number=28
%% source: Q-Net 감정평가사 기출문제 (cst00309) artlSeq=5211789 file=raw/1008/5211789/2023년도 제34회 감정평가사 1차시험 2교시.pdf
%% conversion: kordoc 4.24.51 -> tools/gen1.py (auto); review=proofread-sonnet; flags=none
\begin{question}{28}
공간정보의 구축 및 관리 등에 관한 법령상 용어의 정의에 관한 내용으로 옳지 않은 것은?
\begin{choices}
  \item 자치구가 아닌 구의 구청장은 “지적소관청”이 될 수 있다.
  \item “지목”이란 토지의 주된 용도에 따라 토지의 종류를 구분하여 지적공부에 등록한 것을 말한다.
  \item “경계”란 필지별로 경계점들을 직선으로 연결하여 지적공부에 등록한 선을 말한다.
  \item “등록전환”이란 지적공부에 등록되어 있지 아니한 토지를 지적공부에 등록하는 것을 말한다.
  \item “축척변경”이란 지적도에 등록된 경계점의 정밀도를 높이기 위하여 작은 축척을 큰 축척으로 변경하여 등록하는 것을 말한다.
\end{choices}
\end{question}
\answer{4}
